\let\negmedspace\undefined
\let\negthickspace\undefined
\documentclass[journal]{IEEEtran}
\usepackage[a5paper, margin=10mm, onecolumn]{geometry}
%\usepackage{lmodern} % Ensure lmodern is loaded for pdflatex
\usepackage{tfrupee} % Include tfrupee package

\setlength{\headheight}{1cm} % Set the height of the header box
\setlength{\headsep}{0mm}     % Set the distance between the header box and the top of the text

\usepackage{gvv-book}
\usepackage{gvv}
\usepackage{cite}
\usepackage{physics}
\usepackage{amsmath,amssymb,amsfonts,amsthm}
\usepackage{algorithmic}
\usepackage{graphicx}
\usepackage{textcomp}
\usepackage{xcolor}
\usepackage{txfonts}
\usepackage{listings}
\usepackage{enumitem}
\usepackage{mathtools}
\usepackage{gensymb}
\usepackage{comment}
\usepackage[breaklinks=true]{hyperref}
\usepackage{tkz-euclide} 
\usepackage{listings}
% \usepackage{gvv}                                        
\def\inputGnumericTable{}                                 
\usepackage[latin1]{inputenc}                                
\usepackage{color}                                            
\usepackage{array}                                            
\usepackage{longtable}                                       
\usepackage{calc}                                             
\usepackage{multirow}                                         
\usepackage{hhline}                                           
\usepackage{ifthen}                                           
\usepackage{lscape}
\begin{document}

\bibliographystyle{IEEEtran}

\title{8.4.25}
\author{EE25BTECH11022 - sankeerthan}
% \maketitle
% \newpage
% \bigskip
{\let\newpage\relax\maketitle}

\renewcommand{\thefigure}{\theenumi}
\renewcommand{\thetable}{\theenumi}
\setlength{\intextsep}{10pt} % Space between text and floats


\numberwithin{align}{enumi}
\numberwithin{figure}{enumi}
\renewcommand{\thetable}{\theenumi}

\textbf{Question:}The equation of directrix of parabola $y^2+4y+4x+2=0$
\begin{enumerate}
     \item $x=-1$
     \item $x=\frac{-3}{2}$
     \item $x=1$
     \item $x=\frac{3}{2}$
\end{enumerate}

\textbf{Solution:}
Given the parabola equation:
\begin{align}
y^2 + 4y + 4x + 2 = 0
\end{align}
The equation of a conic with directrix $\vec{n}^\top \vec{x} = c$, eccentricity $e$ and focus $\vec{F}$ is given by
\begin{align}
\vec{x}^\top\vec{V}\vec{x} + 2\vec{u}^\top\vec{x} + f = 0,\vec{x} = \myvec{x \\ y},
\mathbf{V} =\myvec{0 & 0 \\0 & 1},\vec{u} =\myvec{2 \\ 2},f = 2
\end{align}
Find vertex $\vec{c}$ from
\begin{align}
\mathbf{V} \vec{c} &= -\vec{u}  \\
\myvec{0 & 0 \\0 & 1}\myvec{c_x \\c_y} &=-\myvec{2 \\ 2} \\
\myvec{0 \\c_y} &=-\myvec{2 \\ 2} \\
c_y &= -2
\end{align}
Put $c_y=-2$ into parabola to find  $c_x$:
\begin{align}
(-2)^2 + 4(-2) + 4 c_x + 2 = 0  \Rightarrow  c_x = \frac{1}{2}
\vec{c} = \myvec{ \frac{1}{2} \\ -2 }
\end{align}

Eigenvector for zero eigenvalue of $\vec{V}$:
\begin{align}
\vec{a} =\myvec{1 \\ 0}
p = \frac{\|\vec{u}\|^2 - f}{4 \, \vec{u}^T \vec{a}} = \frac{4 + 4 - 2}{8} = \frac{6}{8} = \frac{3}{4}
\end{align}
Equation of directrix (line perpendicular to $\vec{a}$,distance p from vertex):
\begin{align}
\vec{n}^T \vec{x} = \vec{n}^T \vec{c} + p,\text{where}\vec{n}=\vec{a}=\myvec{1\\0}
\vec{x} = \frac{1}{2} + \frac{3}{4} = \frac{5}{4}
\end{align}

Therefore, the equation of the directrix is approximately:
$\vec{x} = \frac{3}{2}$


\begin{figure}
    \centering
    \includegraphics[width=1\columnwidth]{figs/points.png}
    \caption{}
\end{figure}
\end{document}
